
\begin{lemma}\label{lemma:parametrization}
    Let $n = rs$, $K = K_{r,s} \cong S_r \wr S_s \le S_n$.
    There is a bijection between the $(K,S_i\times S_{n-i})$ double cosets of $S_n$ and $P_{r,s}(i)$, the set of partitions of $i$ into at most $s$ parts, all of size at most $r$. 
\end{lemma}

\begin{proof}
    To describe the bijection from $K\backslash S_n/(S_i \times S_{n-i})$ to $P_{r,s}(i)$ explicitly,
    first fix the underlying decomposition $\{1, \ldots, n\} = \{1,\ldots,i\} \cup \{i+1,\ldots, n\}$ for the action of $S_i\times S_{n-i}$. 
    The left cosets in $S_n/(S_i \times S_{n-i})$ are clearly in bijection with the subsets of size $i$ in $\{1, \ldots, n\}$:
    \[
        g(S_i \times S_{n-i}) \longleftrightarrow g(\{1, \ldots, i\}).
    \]
    Now fix a decomposition for the action of $K \cong S_r \wr S_s$: 
    $\{1, \ldots, n\} = B_1 \cup \ldots \cup B_s$,
    where $B_j := \{(j-1)r+1,(j-1)r+2,\ldots,(j-1)r+r\}$ $(j = 1,\ldots,s)$. 
    The elements of $S_s$ permute these blocks,
    and each of the $s$ copies of $S_r$ acts on one of the blocks.
    Given $g \in S_n$, we map the double coset $K g (S_i \times S_{n-i})$ to the partition $\gamma$ which is the non-increasing rearrangement of the sequence
    \[
        (\card{B_1 \cap g(\{1,\ldots,i\})}, \ldots, \card{B_s \cap g(\{1,\ldots,i\})}).
    \]
    This sequence consists of $s$ non-negative integers, each at most $r$, which sum up to $i$. Thus $\gamma \in P_{r,s}(i)$.
    We will show that this map is a bijection.

    For arbitrary $x\in K$ and $y\in S_i\times S_{n-i}$ we have,
    for each $1 \le j \le s$,
    \[
        \card{B_j \cap xgy(\{1,\ldots,i\})} 
        = \card{B_j \cap xg(\{1,\ldots,i\})}
        = \card{x^{-1}(B_j) \cap g(\{1,\ldots,i\})}.
    \]
    The element $x^{-1} \in K = S_r \wr S_s$ permutes the blocks and permutes the elements of each block. This shows that the mapping $Kg(S_i \times S_{n-i}) \mapsto \gamma$ is well defined.

    The mapping from $K\backslash S_n/(S_i \times S_{n-i})$ to $P_{r,s}(i)$ is clearly onto, since
    for each partition $\gamma = (a_1,\ldots,a_s) \in P_{r,s}(i)$ there exists a permutation $g \in S_n$ such that  $\card{B_j \cap g(\{1,\ldots,i\})} = a_j$ $(1 \le j \le s)$. 
		
    Finally, if $K g_1 (S_i \times S_{n-i})$ and $K g_2 (S_i \times S_{n-i})$ are mapped to the same partition~$\gamma$, then there exists a permutation $\pi \in S_s$ satisfying
    \[
        \card{B_j \cap g_1(\{1,\ldots,i\})}
        = \card{B_{\pi(j)} \cap g_2(\{1,\ldots,i\})}
        \qquad (1 \le j \le s).
    \]
    Therefore there exist an element $x \in K = S_r \wr S_s$ such that
    \[
        B_j \cap g_1(\{1,\ldots,i\})
        = B_j \cap xg_2(\{1,\ldots,i\})
        \qquad (1 \le j \le s).
    \]
    It follows that
    \[
        g_1(\{1,\ldots,i\}) = xg_2(\{1,\ldots,i\}),
    \]
    and therefore $g_1 = xg_2y$ for a suitable permutation $y \in S_i \times S_{n-i}$.
\end{proof}

We are now ready to prove Theorem~\ref{prop:e_i_formuli}. 
The proof applies Lemma~\ref{lemma:parametrization} and the explicit bijection described in its proof, combined with Lemma~\ref{lem:char_values_and_inner_product}. 

\begin{proof}[Proof of  Theorem~\ref{prop:e_i_formuli}.]
    Recall that $n=rs$, 
    $K=K_{r,s}\cong S_r\wr S_s\leq S_n$ and 
    $\psi = \phi\uparrow_K^{S_n}$. 
    By Frobenius reciprocity (twice) and 
    Mackey's formula~\cite[(5.2), Problem (5.6)]{Isaacs},
    \begin{equation}\label{Mackey_sum}
    \begin{split}
        e_i 
        &= \langle \psi, \chi^{(1^i) \oplus (n-i)} \rangle 
        = \langle \phi\uparrow_K^{S_n}, (\varepsilon_{S_i} \times 1_{S_{n-i}})\uparrow_{S_i \times S_{n-i}}^{S_n} \rangle \\
        &= \langle {\phi\uparrow_K^{S_n}}\downarrow_{S_i \times S_{n-i}}^{S_n}, \varepsilon_{S_i}\times 1_{S_{n-i}} \rangle \\
        &= 
        \sum_{[g] \in K \backslash S_n / (S_i \times S_{n-i})}
        \langle \phi^g\downarrow^{K^g}_{K^g \cap (S_i \times S_{n-i})}\uparrow_{K^g \cap (S_i \times S_{n-i})}^{S_i \times S_{n-i}}, \varepsilon_{S_i} \times 1_{S_{n-i}} \rangle \\
        &=
        \sum_{[g] \in K \backslash S_n / (S_i \times S_{n-i})}
        \langle \phi^g\downarrow^{K^g}_{K^g \cap (S_i \times S_{n-i})}, (\varepsilon_{S_i} \times 1_{S_{n-i}})\downarrow^{S_i \times S_{n-i}}_{K^g \cap  (S_i \times S_{n-i})} \rangle.
    \end{split}
    \end{equation}
		
    The above sums are indexed by the $(K,S_i\times S_{n-i})$ double cosets in $S_n$. 
    For each representative $g$ of a double coset, 
    $K^g := g^{-1}Kg$ is the corresponding conjugate of $K$ 
    and $\phi^g$ is the character on $K^g$ defined by
    $\phi^g(g^{-1}kg) := \phi(k)$ for all $k \in K$.

    By Lemma~\ref{lemma:parametrization}, 
    these double cosets are parametrized by the partitions in $P_{r,s}(i)$.  Let us determine the summand of \eqref{Mackey_sum} 
    corresponding to a partition $\gamma\in P_{r,s}(i)$ in which part $j$ occurs with multiplicity $k_j = k_j(\gamma)$ $(0 \le j \le r)$. 
    By the bijection described in the proof of Lemma~\ref{lemma:parametrization}, 
    \[
        k_j = \card{\{t \,:\, \card{B_t \cap g(\{1, \ldots, i\})} = j\}}
        \qquad (0 \le j \le r),
    \]
    where $B_t := \{(t-1)r+1,\ldots,(t-1)r+r\}$ $(1 \le t \le s)$. Thus
    \[
        K^g \cap (S_i \times S_{n-i})
        \cong (R_{r,0} \wr S_{k_0}) \times (R_{r,1} \wr S_{k_1}) \times \cdots \times (R_{r,r}\wr S_{k_r}),
    \]
    where $R_{r,j} = S_j \times S_{r-j}$, as above. In particular,
    \begin{align*}
        (\varepsilon_{S_i} \times 1_{S_{n-i}})\downarrow^{S_i \times S_{n-i}}_{K^g \cap (S_i \times S_{n-i})}
        &= \bigotimes_j (\varepsilon_{S_{jk_j}} \times 1_{S_{(r-j)k_j}})\downarrow^{S_{jk_j} \times S_{(r-j)k_j}}_{R_{r,j}\wr S_{k_j}}
        = \bigotimes_j \nu_{r,j,k_j}.
    \end{align*}
    Note that, by Observation~\ref{obs:phi_product}, $\phi_{r,s}$ factors similarly.
    Therefore the corresponding summand in~\eqref{Mackey_sum} is
    \begin{equation}\label{Mackey_summand}
        \langle \phi^g\downarrow^{K^g}_{K^g \cap (S_i \times S_{n-i})},
        (\varepsilon_{S_i} \times 1_{S_{n-i}})\downarrow^{S_i \times S_{n-i}}_{K^g \cap (S_i \times S_{n-i})} \rangle 
        = \prod_{j=0}^{r} \langle \phi_{r,k_j}\downarrow^{S_r\wr S_{k_j}}_{R_{r,j}\wr S_{k_j}}, \nu_{r,j,k_j} \rangle.
    \end{equation}
		
    We have already computed these inner products in Lemma~\ref{lem:char_values_and_inner_product}.
    By \eqref{Mackey_sum}, \eqref{Mackey_summand}, Lemma~\ref{lem:char_values_and_inner_product}
    and Definition~\ref{def:f} we obtain
    \[
        e_{i,(r^s)} 
        = \langle \phi\uparrow_K^{S_n}, (\varepsilon_{S_i} \times 1_{S_{n-i}})\uparrow_{S_i \times S_{n-i}}^{S_n} \rangle
        = \sum_{\gamma\in P_{r,s}(i)} \prod_{j=0}^{r} 
        (-1)^{(j+1)k_j(\gamma)} \binom{(-1)^{j+1} f_j(r)}{k_j(\gamma)},
    \]
    as claimed.
    If $s=1$ then $P_{r,1}(i)$ contains (for $0 \le i \le r$) 
    a unique partition $\gamma = (i)$, for which $k_i(\gamma) = 1$ is the unique nonzero multiplicity. Thus, in this case,  
    \[
        e_{i,(r)} 
        = (-1)^{i+1} \binom{(-1)^{i+1}f_i(r)}{1} 
        = f_i(r).
        \qedhere
    \]
\end{proof}
 
	
\subsection{A product formula}
\label{sec:prod}
	
Now we derive a restatement of Theorem~\ref{prop:e_i_formuli} as a product formula for formal power series. 

\begin{corollary}\label{power_series_form}
    For a positive integer $r$, define the formal power series
    \[
        E_r(x,y) := \sum_{i,s \ge 0} e_{i,(r^s)} x^i y^s.
    \]
    Then
    \[
        E_r(x,y) 
        = \prod_{j=0}^{r} (1 - (-x)^j y)^{(-1)^{j+1}f_j(r)}.
    \]
\end{corollary}
	
\begin{proof}
    Recall the following formal power series expansion, valid for any integer $f$:
    \[
        (1+t)^{f} = \sum_{n=0}^\infty \binom{f}{n} t^n.
    \]
    By Theorem~\ref{prop:e_i_formuli},
    with the obvious extension for $s=0$,
    \begin{align*}
        E_r(x,y) 
        &= \sum_{i,s \ge 0} \left(
        \sum_{\substack{k_0,\ldots,k_r \ge 0 \\ \sum_j k_j=s \\ \sum_j jk_j = i}} 
        \prod_{j = 0}^{r}
        (-1)^{(j+1)k_j} \binom{(-1)^{j+1} f_j(r)}{k_j}
        \right) x^i y^s \\
        &= \prod_{j=0}^{r} \sum_{k_j=0}^\infty
        (-1)^{(j+1)k_j} \binom{(-1)^{j+1} f_j(r)}{k_j}
        x^{jk_j} y^{k_j} \\
        &= \prod_{j=0}^{r} (1 + (-1)^{j+1} x^j y)^{(-1)^{j+1}f_j(r)},
    \end{align*}
    as required.
\end{proof}
	
\begin{remark}\label{rem:le_3}
    For small $r$ and arbitrary $s$, Corollary~\ref{power_series_form}  enables us to determine explicitly the hook-multiplicities $m_{i,(r^s)}$. 
    This is done by recalling Equation~\eqref{eq:alternating} and the fact that, by definition, $e_{i,(r^s)}$ is the coefficient of $x^i y^s$ in $E_r(x,y)$. 
    For example, by Corollary~\ref{power_series_form} and Observation~\ref{obs:f_values}, 
    $E_2(x,y) = E_3(x,y) = (1+xy)(1-x^2y)^{-1}$. 
    Thus, for $r \in \{2,3\}$ and any $s \ge 1$, the value of $e_{i,(r^s)}$ is $1$ for $i \in \{2s-1,2s\}$ and zero otherwise. 
    Combining this with Equation~\eqref{eq:alternating}, 
    it follows that, for $r \in \{2,3\}$ and $s \ge 1$, the hook multiplicity $m_{i,(r^s)}$ 
    is 1 for $i = 2s-1$ and zero otherwise. 
\end{remark}	


\subsection{Non-negativity}\label{sec:nn}
		
Now we are ready to prove Theorem~\ref{t:main_hook-alternating2}.
	
\begin{proof}[Proof of Theorem~\ref{t:main_hook-alternating2}.]
    Assume that $r$ is not square-free.
	Recall, from Definition~\ref{def:F_polynomial}, the notation $F_r(x) := \sum_{j=0}^{r} f_j(r) x^j$.
	By 	Proposition~\ref{prop:cycles_nonnegativity} and Corollary~\ref{t:F_and_M}	
    we may write $F_r(x) = (1+x)^2 G_r(x)$, where $G_r(x) = \sum_{j=0}^{r-2} g_j(r) x^j$ is a polynomial with non-negative integer coefficients. 
	Let $g_j(r) := 0$ for $j < 0$ or $j > r-2$. Then
	\[
	   f_j(r) = g_j(r) + 2g_{j-1}(r) + g_{j-2}(r)
	   \qquad (\forall j).
	\]
    Therefore, by Corollary~\ref{power_series_form},
    \begin{equation*}\begin{split}
        E_r(x,y) 
        &= \prod_{j \ge 0} (1-(-x)^j y)^{(-1)^{j+1}f_j(r)} \\
        &= \prod_{j \ge 0} (1-(-x)^j y)^{(-1)^{j+1}(g_j(r) + 2g_{j-1}(r) + g_{j-2}(r))} \\
        &= \prod_{j \ge 0} \left(
        \frac{(1 - (-x)^{j} y)(1 - (-x)^{j+2} y)}{(1 - (-x)^{j+1} y)^2} \right)^{(-1)^{j+1} g_j(r)}.
    \end{split}
    \end{equation*}
    We claim that each factor in 
    this product has the form $1 + (1+x)^2 p_j(x,y)$, with $p_j(x,y)$ a formal power series with non-negative integer coefficients. 
    This implies that $(E_r(x,y)-1)/(1+x)^2$ is itself a formal power series with non-negative integer coefficients, completing the proof of Theorem~\ref{t:main_hook-alternating2}.
	
    Indeed, if $j$ is odd then the corresponding factor is
    \[
        \left( \frac{(1 + x^{j} y)(1 + x^{j+2} y)}{(1 - x^{j+1} y)^2} \right)^{g_j(r)}
        = \left( 1 + \frac{(x+1)^2 x^j y}{(1 - x^{j+1} y)^2} \right)^{g_j(r)},
    \]
    where
    \[
        \frac{x^j y}{(1 - x^{j+1} y)^2}
        = x^j y \cdot \sum_{i \ge 0} (i+1) (x^{j+1} y)^i
    \]
    is a formal power series with non-negative integer coefficients.
	
    Finally, if $j$ is even then the corresponding factor is
    \[
        \left( \frac{(1 + x^{j+1} y)^2}{(1 - x^{j} y)(1 - x^{j+2} y)} \right)^{g_j(r)}
        = \left( 1 + \frac{(x+1)^2 x^j y}{(1 - x^{j} y)(1 - x^{j+2} y)} \right)^{g_j(r)},
    \]
    where
    \[
        \frac{x^j y}{(1 - x^{j} y)(1 - x^{j+2} y)}
        = x^j y \cdot \sum_{i \ge 0} (x^{j} y)^i \cdot \sum_{k \ge 0} (x^{j+2} y)^k
    \]
    is a formal power series with non-negative integer coefficients.
\end{proof}


\section{Additional results}
\label{sec:final}

\subsection{Combinatorial identities}
\label{remarks_on_s=1}

In this subsection it will be shown that Lemma~\ref{lem:char_eval_cycles} and Theorem~\ref{prop:e_i_formuli} imply well-known combinatorial identities.

Recall the {\em major index} of a permutation $\pi \in S_n$,
\[
    \maj(\pi):=\sum\limits_{i\in \Des(\pi)}i.
\]
The following identity is due to 
Garsia~\cite{Garsia}. A purely combinatorial proof was given by Wachs~\cite{Wachs}.

\begin{proposition}[{\cite[Equation 5.8]{Garsia}}]\label{prop:Garsia}
     For every partition $\lambda\vdash n$,
    \[
        \sum_{\pi \in \C_\lambda} \zeta^{\maj(\pi)}
        = \begin{cases} 
            \mu(r), &\text{if } \lambda = (r^s); \\
    		0, &\text{otherwise },
		\end{cases}
    \]
    where $\zeta$ is a primitive $n$-th root of unity and $\mu$ is the M\"obius function. 
\end{proposition}

The following lemma follows from the work of Stembridge~\cite{Stembridge}.

\begin{lemma}\label{prop:equiv_Garsia}
    For every Schur-positive set $\A\subseteq S_n$ with associated $S_n$-character $\phi := \ch^{-1}(\Q(\A))$, the value of $\phi$ at an $n$-cycle $c\in S_n$ is
    \[
        \phi(c) = \sum_{\pi \in \A} \zeta^{\maj(\pi)},
    \]
    where $\zeta$ is a primitive $n$-th root of unity. 
\end{lemma}

\begin{proof}
    First, recall 
    the definition of the descent set of standard Young tableaux (SYT) 
    from Equation~\eqref{def:Des_SYT}. By~\cite[Lemma 3.4]{Stembridge}, for every partition $\nu\vdash n$ the value of the irreducible $S_n$-character $\chi^\nu$ at an $n$-cycle $c\in S_n$ is 
    \begin{equation}\label{eq:Stembridge}
        \chi^\nu(c)=\sum\limits_{T\in \SYT(\nu)}\zeta^{\maj(T)}.
    \end{equation}
    Let $\A\subseteq S_n$ be Schur-positive with associated $S_n$-character $\phi$, i.e., 
    $\Q(A)=\ch (\phi)$. 
    By Lemma~\ref{lem:Schur-gf} 
    together with Equation~\eqref{eq:Stembridge},
    \[
        \sum_{\pi \in \A} \zeta^{\maj(\pi)}
        = \sum_{\nu \vdash n} \langle \Q(\A), s_\nu \rangle \sum_{T \in \SYT(\nu)} \zeta^{\maj(T)}
        = \sum_{\nu \vdash n} \langle \phi, \chi^\nu \rangle \chi^\nu(c)
        = \phi(c), 
    \]
    completing the proof.
\end{proof}

In light of Lemma~\ref{prop:equiv_Garsia} we deduce the following. 

\begin{corollary}
    Lemma~\ref{lem:char_eval_cycles} is equivalent to Garsia's identity (Proposition~\ref{prop:Garsia}).
\end{corollary}

\begin{proof}
    By the Gessel--Reutenauer Theorem (Theorem~\ref{thm:GR}), for every $\lambda \vdash n$, 
    the conjugacy class of cycle type $\lambda$ is Schur-positive with $\ch^{-1}(\Q(\C_\lambda))=\psi^\lambda$. 
    Letting $\A = \C_\lambda$ in Lemma~\ref{prop:equiv_Garsia}, 
    Proposition~\ref{prop:Garsia} implies Lemma~\ref{lem:char_eval_cycles} and vice versa.
\end{proof}

There is a combinatorial description, due to Schocker,
of the multiplicity of an arbitrary irreducible character of $S_n$ in the higher Lie character. 
In its full generality it is too complicated to be presented here, see~\cite{schocker} for the details. 
The special case of a full cycle, $\lambda=(n)$ (for which
$\psi^{(n)} = \omega^{(n)}\uparrow_{\ZZ_n}^{S_n}$ is the Lie character), is due 
to Kra\'skiewicz and Weyman~\cite{KraskiewiczWeyman}. 
	
\begin{theorem}[{Kra\'skiewicz--Weyman, see~\cite[Theorem 8.4]{Garsia}}]\label{thm:KW}
    For every partition \mbox{$\nu\vdash n$}, the multiplicity $m_{\nu,(n)} := \langle \psi^{(n)},\chi^{\nu} \rangle$ is equal to the cardinality of the set
    \[
        \{T \in \SYT(\nu):\ \maj(T) \equiv 1 \pmod n\}.
    \]	
\end{theorem}
	
\begin{corollary}\label{lem:Kra-Wey}
    For every $0\leq k\leq n$, the multiplicity 
    $m_{k,(n)} := \langle \psi^{(n)},\chi^{(n-k,1^k)} \rangle$ 
    is equal to the cardinality of the set
    \[
        \left\{ 1 \le a_1 < \cdots < a_k \le n-1 : \sum_{i=1}^k a_i \equiv 1 \!\!\!\!\pmod{n} \right\}.
    \]
\end{corollary}
	
\begin{proof}
    The map $\Des: \SYT(n-k,1^k) \rightarrow \binom{[n-1]}{k}$ is a bijection, where $\binom{[n-1]}{k}$ denotes the set of all $k$-subsets of $[n-1]$.
    The major index of a tableau is the sum of the elements of its descent set.
\end{proof}
	
Consider the following combinatorial identity.

\begin{proposition}\label{t:equivalence_KW}
    For every $0 \le k \le n$,
    \[
        \left| \left\{ 1 \le a_1 < \cdots < a_k \le n : \sum_{i=1}^k a_i\equiv 1 \!\!\!\!\pmod{n} \right\} \right| 
        = \sum_{d|(n,k)} \frac{\mu(d)(-1)^{(d+1)k/d}}{n} \binom{n/d}{k/d}.
    \]
\end{proposition}
	
\begin{proof}
    Let 
    $H(x,q):=\prod_{i=1}^n(1+xq^i)$. 
    Writing
    \[
        H(x,q) \equiv \sum_{k=0}^{n} \sum_{t=0}^{n-1} c_{k,t} x^kq^t \mod (q^n-1),
    \] 
    the cardinality we are interested in is the coefficient $c_{k,1}$ of $x^kq$.
    For any $d|n$ and $\eta$ a primitive $d$-th root of unity (we write $o(\eta)= d$), 
    \[
        H(x,\eta) 
        = \prod_{i=1}^n (1+x\eta^i) 
        = \left( \prod_{i=1}^d (1+x\eta^i) \right)^{n/d} 
        = (1-(-x)^d)^{n/d},
    \]
    where the last equality holds since both sides are polynomials of the same degree, with exactly the same roots and the same constant term. 
    
    Let $\omega$ be a primitive $n$-th root of unity. Then
    \[
        H(x,\omega^j) = \sum_{t=0}^{n-1} h_t(x)\omega^{jt}
        \qquad (0 \le j \le n-1),
    \]
    where $h_t(x) = \sum_{k=0}^n c_{k,t} x^k$. 
    Fourier inversion gives
    \begin{align*}
        \sum_{k=0}^{n} c_{k,1}x^k 
        &= h_1(x)
        = \frac{1}{n} \sum_{j=0}^{n-1} H(x,\omega^j) \omega^{-j} \\
        &= \frac{1}{n} \sum_{d|n}  (1-(-x)^d)^{n/d}
        \sum_{j \,:\, o(\omega^j)=d} \omega^{-j} \\
        &= \frac{1}{n} \sum_{d|n}  (1-(-x)^d)^{n/d}
        \mu(d) \\
        &= \frac{1}{n} \sum_{k=0}^{n} \sum_{d|(n,k)} \binom{n/d}{k/d} (-1)^{(d+1)k/d} x^k \mu(d),
    \end{align*}
    as required.
\end{proof}

\begin{observation}\label{prop:equiv_f_KW}
    Proposition~\ref{prop:e_i_formuli1} is equivalent to Proposition~\ref{t:equivalence_KW}.
\end{observation}

\begin{proof}
    By considering separately the cases $a_k < n$ and $a_k = n$, Corollary~\ref{lem:Kra-Wey} yields 
    \[
        \left| \left\{ 1 \le a_1 < \cdots < a_k \le n : \sum_{i=1}^k a_i \equiv 1 \!\!\!\!\pmod{n} \right\} \right|
        = m_{k,(n)} + m_{k-1,(n)} 
        = e_{k,(n)}.
    \]
    By Proposition~\ref{prop:e_i_formuli1} and Definition~\ref{def:f},
    \[
        e_{k,(n)} 
        = f_k(n)
        = \sum_{d|(n,k)} \frac{\mu(d)(-1)^{(d+1)k/d}}{n} \binom{n/d}{k/d},
    \]
    proving Proposition~\ref{t:equivalence_KW}.
    The opposite direction is similar. 
\end{proof}
    

\begin{remark}
    Noting that Proposition~\ref{prop:e_i_formuli1} is the special case $s=1$ of Theorem~\ref{prop:e_i_formuli}, one concludes that Proposition~\ref{t:equivalence_KW} is a consequence of the latter.
\end{remark}

It remains a challenge to find such a direct link between Schocker's general description of the multiplicity and our version in Theorem~\ref{prop:e_i_formuli}.

\begin{remark}\label{rem:ET}
    Proposition~\ref{prop:e_i_formuli1} 
    is not new.
    For example, by the Gessel--Reutenauer Theorem (Theorem~\ref{thm:GR}) together with Observation~\ref{lem:1}, 
    Theorem~\ref{prop:e_i_formuli} at $s=1$ is equivalent to the equation
    \[
        |\{\pi\in \C_{(n)}: \ \Des(\pi)=[j]\}|
        = \frac{1}{n} \sum_{d|n} \mu(d) (-1)^{j - \lfloor j/d \rfloor} \binom{n-1}{j-1}
        \qquad (0 \le j <n),        
    \]
    which is an immediate consequence of a recent result of Elizalde and Troyka~\cite[ Theorem 3.1]{ET}.  
    An older proof was presented to us by Sheila Sundaram~\cite{Sundaram_personal}, deducing 
    Proposition~\ref{prop:e_i_formuli1} 
    from~\cite[Lemma 2.7]{Sundaram_Adv}. 
    The reader is referred to~\cite{ET} for further discussion and relations to the enumeration of Lyndon words.  
\end{remark}


\subsection{Cellini's cyclic descents}\label{sec:Cellini}
	
In this subsection it is shown that the apparently natural approach does not provide a cyclic descent extension for many conjugacy classes in $S_n$.
	
Recall the original notion of cyclic descent set defined by Cellini~\cite{Cellini},
\[
	\CDes(\pi) 
	:= \{1 \leq i \leq n \,:\, \pi_i > \pi_{i+1} \}
    \qquad (\forall \pi \in S_n),
\]
with the convention $\pi_{n+1}:=\pi_1$.
	
Elizalde and Roichman~\cite{ER19} presented several subsets of $S_n$ on which the image of Cellini's cyclic descent map is closed under cyclic rotation, thus leading to a cyclic extension of $\Des$.
However, as we shall see, many conjugacy classes do not have this property. In fact, we conjecture that only two conjugacy classes have this property.
Here are some partial results.
	
\begin{proposition}\label{prop:Cel1}
	For $n = rs > 1$, the image of Cellini's cyclic descent map on any conjugacy class of cycle type $(r^s)$ is not closed under cyclic rotation.
\end{proposition}
	
\begin{proof}
	Recall the notation $\C_\lambda$ from Section~\ref{sec:background}.
	For $r=1$ and $n = s > 1$, $\C_{(1^n)}$ consists of the identity permutation only. Cellini's $\CDes(id)$ is the singleton $\{n\}$ and, for~$n > 1$, the set $\{\{n\}\}$ is not closed under cyclic rotation.
	For $r>1$ (and $s \ge 1$), 
    let~$\sigma = [s+1,s+2,\dots,n,1,2,\dots,s]$.
    In other words, $\sigma$ is the permutation in $S_n$ defined by 
	\[
		\sigma(i) = i+s \pmod{n} \qquad (\forall i\in [n]).
	\]
	Then $\sigma \in \C_{(r^s)}$, 
    and Cellini's $\CDes(\sigma)$ is the singleton $\{n-s\}$. 
    If the image of $\CDes$ on $\C_{(r^s)}$ is invariant under cyclic rotation then there is a permutation $\pi \in \C_{(r^s)}$ with $\CDes(\pi) = \{n\}$. The only permutation in $S_n$ with this property is the identity permutation, which is not in $\C_{(r^s)}$. This is a contradiction.
\end{proof}
	
\begin{proposition}
	For $n > 1$, the image of Cellini's cyclic descent map on any conjugacy class of $k$-cycles in $S_n$, except $2$-cycles in $S_3$ and $3$-cycles in $S_4$, is not closed under cyclic rotation.
\end{proposition}
	
\begin{proof}
	Letting $r=n$ in Proposition~\ref{prop:Cel1}, statement holds on $n$-cycles.
	For $k<n$ let $\sigma\in \C_{(k,1^{n-k})}$ be the permutation $[k,1,2,\dots,k-1,k+1,k+2,\dots,n]$. Then $\CDes(\sigma)=\{1,n\}$.
	By the equivariance property,  there must be a $k$-cycle $\pi$ with cyclic descent set $\{1,2\}$. Then $\pi(3)$ is the minimal value thus it is equal to 1, and $\pi(1)$
	is the maximal value thus it is equal to $n$. Let $\pi(2)=x$, thus
	\[
		\pi=[n,x,1,2,\dots,x-1,x+1,\dots,n-1],
	\]
	namely, for every $3<i\le x+1$, $\pi(i)=i-2$ and for every $x+1<i\le n$, $\pi(i)=i-1$. Thus $\pi$ has no fixed points unless $x=2$, in this case $\pi$
	has cycle type $(n-1,1)$.  One deduces that statement holds for $k<n-1$.
		
	For the $(n-1)$-cycles an argument similar to the above works. For $n>4$ the permutation $\sigma=(1,n-1,n-2,\ldots,5,4,2,3)(n)$, i.e., $[n-1,3,1,2,4,5,\ldots,n-2,n]\in \C_{(n-1,1)}$ has cyclic descent set $\CDes(\sigma)=\{1,2,n\}$. 
	By the equivariance property, there is a permutation $\pi\in \C_{(n-1,1)}$ with cyclic descent set $\{1,n-1,n\}$. Then $\pi(2)$ is the minimal value thus it is equal to $1$, and $\pi(n-1)$ is the maximal value thus it is equal to $n$. If $\pi(n)=l$ and $\pi(1)=k$ then $1<k<l<n$ and $\pi=[k,1,\dots,k-1,k+1,...,l-1,l+1,\dots,n,l]$, 
	so the cycle $(1,k,k-1,\ldots,2)$ of $\pi$ has length $2\le k<n-1$,  
	contradicting the equivariance property.
\end{proof}

\begin{conjecture} 
    For $n > 1$, the image of Cellini's cyclic descent map on a conjugacy class $\C$ is invariant under cyclic rotation if and only if $n \in \{3,4\}$ and $\C$ is the conjugacy class of $(n-1)$-cycles.
\end{conjecture}


\subsection{Palindromicity of hook multiplicities}
\label{sec:palindrom}	

A sequence $a_0,\ldots,a_n$ (equivalently, the polynomial $a_0+a_1x+\cdots+a_nx^n$) is 	
{\em palindromic} (or {\em symmetric})
if $a_i=a_{n-i}$ for all $0 \le i \le n$. 

For a partition $\lambda\vdash n$ recall the notation
\[
    m_{k,\lambda}
    := \langle \psi^{\lambda},\chi^{(n-k,1^k)} \rangle 
    \qquad (0\le k<n).
\]

In this subsection we prove the following. 

\begin{proposition}\label{conj:r2}
    Consider the partition $\lambda = (r^s)$ for positive integers $r$ and $s$.
    \begin{enumerate}
        \item 
        If $s=1$ then the hook-multiplicity sequence  $m_{0,(r)},\,m_{1,(r)},\ldots, m_{r-1,(r)}$ is palindromic if and only if either $r$ is odd or $r\equiv 0 \pmod 4$. 

        \item
        If $s>1$ then the hook-multiplicity sequence $m_{0,(r^s)},\,m_{1,(r^s)},\ldots, m_{rs-1,(r^s)}$ is palindromic if and only if $r\equiv 0 \pmod 4$.
    \end{enumerate}
\end{proposition}

\begin{proof}

{\bf 1.} 
Assume first that $s=1$.
Recall the notations $M_{(r)}(x) := \sum_{j=0}^{r-1} m_{j,(r)}x^j$ and 
$F_r(x) := \sum_{j=0}^{r} f_{j}(r)x^j$.  
By Corollary~\ref{t:F_and_M},
$F_r(x) = (1+x)M_{(r)}(x)$. 
Hence $M_{(r)}(x)$ is palindromic if and only if $F_r(x)$ is palindromic.  

If $r$ is odd then, for every $j$, every divisor $d|(r,j)$ is odd, hence $(d+1)j/d$ 
is even. 
Thus, by Definition~\ref{def:f}, 
\[
    f_j(r)
    = \sum_{d|(r,j)} \frac{\mu(d)}{r} \binom{r/d}{j/d} 
    = \sum_{d|(r,r-j)} \frac{\mu(d)}{r} \binom{r/d}{(r-j)/d}
    = f_{r-j}(r) 
    \qquad (0 \le j \le r),
\]
and $F_r(x)$ is palindromic.
 
Next consider the case $r\equiv 0 \pmod 4$. 
Since $\mu(d)=0$ for $d\equiv 0 \pmod 4$, Definition~\ref{def:f} implies that  
\begin{align*}
    f_j(r)
    &= \sum_{d|(r,j)} 
    \frac{\mu(d)(-1)^{(d+1)j/d}}{r} \binom{r/d}{j/d} \\
    &= \sum_{\substack{d|(r,j) \\ d \text{ odd}}}
    \frac{\mu(d)(-1)^{(d+1)j/d}}{r} \binom{r/d}{j/d}
    + \sum_{\substack{d|(r,j) \\ d \equiv 2\!\!\!\! \pmod 4}}
    \frac{\mu(d)(-1)^{(d+1)j/d}}{r} \binom{r/d}{j/d}.
\end{align*}
Again, if $d|(r,j)$ is odd then $(d+1)j/d$ is even. 
If $d \equiv 2 \pmod 4$ then
\[
    (-1)^{(d+1)j/d} (-1)^{(d+1)(r-j)/d}
    = (-1)^{(d+1)r/d} = 1,
\]
since $r/d$ is even. 
It follows that
\[
    (-1)^{(d+1)j/d} = (-1)^{(d+1)(r-j)/d}
\]
in this case. 
We deduce that if $r\equiv 0\pmod 4$ then, for every $0 \le j \le r$,
\begin{align*}
    f_j(r) 
    &= \sum_{\substack{d|(r,j) \\ d \text{ odd}}} 
    \frac{\mu(d)}{r} \binom{r/d}{j/d} 
    + \sum_{\substack{d|(r,j) \\ d \equiv 2\!\!\!\! \pmod 4}}
    \frac{\mu(d)(-1)^{(d+1)j/d}}{r} \binom{r/d}{j/d} \\
    &= \sum_{\substack{d|(r,r-j) \\ d \text{ odd}}}
    \frac{\mu(d)}{r} \binom{r/d}{(r-j)/d}
    + \sum_{\substack{d|(r,r-j) \\ d \equiv 2\!\!\!\! \pmod 4}}
    \frac{\mu(d)(-1)^{(d+1)(r-j)/d}}{r} \binom{r/d}{(r-j)/d} \\
    &= f_{r-j}(r),  
\end{align*}
hence $F_r(x)$ is palindromic.

On the other hand, if 
$r\equiv 2 \pmod 4$ then,  
letting $j = 2$,
\[
    f_2(r) 
    = \frac{1}{r} \left[ \binom{r}{2} + \binom{r/2}{1} \right]
    = \frac{r}{2}
    \ne \frac{r-2}{2}
    = \frac{1}{r} \left[ \binom{r}{r-2} - \binom{r/2}{(r-2)/2} \right]
    = f_{r-2}(r),
\]
thus $F_r(x)$ is not palindromic in this case.

{\bf 2.}  
Assume now that $s > 1$.
By Equation~\eqref{eq:alternating}, 
\[
    \sum_{j=0}^{rs} e_{j,(r^s)}x^j
    = (1+x) \sum_{j=0}^{rs-1} m_{j,(r^s)}x^j.
\]
Thus the sequence $m_{0,(r^s)}, \ldots m_{rs-1,(r^s)}$ is palindromic if and only if the sequence $e_{0,(r^s)}, \ldots, e_{rs,(r^s)}$ is.
We will show that this happens if and only if $r \equiv 0 \pmod 4$.

Assume first that $r \equiv 0 \pmod 4$.
Following Definition~\ref{def:partition},
for each partition 
$\gamma  = (\gamma_1, \ldots, \gamma_s) \in P_{r,s}(i)$ 
consider the complementary partition 
$\bgamma = (\bgamma_1, \ldots, \bgamma_s) \in P_{r,s}(rs-i)$, defined by  
$\bgamma_\ell := r-\gamma_{s+1-\ell}$ $(1 \le \ell \le s)$,
or equivalently by
$k_j(\bgamma) = k_{r-j}(\gamma)$ $(0 \le j \le r)$.
Since $r \equiv 0 \pmod 4$,
$f_j(r) = f_{r-j}(r)$ $(0 \le j \le r)$,
by Part 1 of the current proof.
In addition, $j+1$ and $r-j+1$ have the same parity when $r$ is even and $j$ is arbitrary. 
Using Proposition~\ref{prop:e_i_formuli},
it follows that for $r\equiv 0 \pmod 4$ and any $0 \le i \le rs$,
\begin{align*}
    e_{i,(r^s)}
    &= \sum_{\gamma \in P_{r,s}(i)} \prod_{j=0}^{r}
    (-1)^{(j+1)k_j(\gamma)} \binom{(-1)^{j+1} f_j(r)}{k_j(\gamma)} \\
    &= \sum_{\gamma \in P_{r,s}(rs-i)} \prod_{j=0}^{r}
    (-1)^{(r-j+1)k_{r-j}(\gamma)} \binom{(-1)^{r-j+1} f_{r-j}(r)}{k_{r-j}(\gamma)}= e_{rs-i,(r^s)},  
\end{align*}
proving palindromicity in this case. 		
For the converse,
    consider again the explicit formula for $e_{i,(r^s)}$ (from Proposition~\ref{prop:e_i_formuli}) written above. 
    Each summand corresponds to a partition $\gamma \in P_{r,s}(i)$. 
    According to Remark~\ref{rem:when_factor_is_zero}, the summand is zero if and only if either
    $k_j(\gamma) > f_j(r)$ for some odd $j$, or
    $k_j(\gamma) > 0 = f_j(r)$ for some even $j$.
    
       We have
        $f_0(r) = 0$ for $r>1$, 
        $f_r(r) = 0$ for $r>2$, and
        $f_1(r) = f_{r-1}(r) = 1$ for $r>0$ thanks to Observation~\ref{obs:f_values}.
        It follows that for $\gamma \in P_{r,s}(i)$ to contribute a nonzero summand, it is necessary that
        $k_0(\gamma) = 0$ (for $r>1$), 
        $k_r(\gamma) = 0$ (for $r>2$), 
        $k_1(\gamma) \in \{0,1\}$ (for $r>0$) and 
        $k_{r-1}(\gamma) \in \{0,1\}$ (for $r-1$ odd).
        
        Assume first that $r>1$ is odd. 
        The restrictions on $k_0(\gamma)$ and $k_1(\gamma)$ imply that for~$i=s$ there is no relevant $\gamma \in P_{r,s}(s)$
        (since $s > 1$). 
        The restriction on $k_r(\gamma)$ implies that for $i=rs-s$ there is only one relevant $\gamma \in P_{r,s}(rs-s)$, with $k_{r-1}(\gamma) = s$ (and $k_j(\gamma) = 0$ for all other values of $j$).
        Thus 
        \[
            e_{rs-s,(r^s)}
            = \binom{s}{s}
            = 1 > 0
            = e_{s,(r^s)},
        \]
        and there is no palindromicity.
        
        If $r=1$ then $f_0(1) = f_1(1) = 1$ and the unique partition $\gamma \in P_{1,s}(i)$ has $k_1(\gamma) = i$ and $k_0(\gamma) = s-i$ $(0 \le i \le s)$. It follows that
        \[
            e_{i,(1^s)} 
            = \binom{s-i}{s-i} \binom{1}{i}
            = \begin{cases}
                1, &\text{if } i \in \{0,1\}; \\
                0, &\text{otherwise.}
            \end{cases}
        \]
        For $s>1$ this sequence is not palindromic.
        
        Assume now that $2 < r \equiv 2 \pmod 4$.
        The restrictions on $k_0(\gamma)$ and $k_1(\gamma)$ imply (actually, for any $r>1$) that 
        $e_{i,(r^s)}=0$ for $0 \le i < 2s-1$ and 
        $e_{2s-1,(r^s)} = \binom{f_2(r)+s-2}{s-1}$.
        Similarly, the restrictions on $k_r(\gamma)$ and $k_{r-1}(\gamma)$ imply (for even $r>2$) that 
        $e_{rs-i,(r^s)}=0$ for $0 \le i < 2s-1$ and 
        $e_{rs-2s+1,(r^s)} = \binom{f_{r-2}(r)+s-2}{s-1}$.  
Recall, from Part 1 of the current proof, that for $r \equiv 2 \pmod 4$
\[
    f_2(r) 
    = \frac{r}{2} 
    > \frac{r-2}{2} 
    = f_{r-2}(r).
\]
Therefore
\[
    e_{2s-1,(r^s)}
    = \binom{f_2(r)+s-2}{s-1} 
    > \binom{f_{r-2}(r)+s-2}{s-1}
    = e_{rs-2s+1,(r^s)}
\]
and the sequence is not palindromic.

Finally, if $r=2$ then,
by Remark~\ref{rem:le_3},
\[
    e_{i,(2^s)} 
    = \begin{cases}
        1, &\text{if } i \in \{2s-1,2s\}; \\
        0, &\text{otherwise}
    \end{cases}
\]
and this sequence is not palindromic.
This completes the proof. 
\end{proof}


\section{Final remarks and open problems}\label{sec:open}

Recall the notation $m_{k,\lambda}:= \langle 
\psi^{\lambda},\chi^{(n-k,1^k)} \rangle$. 
By Proposition~\ref{prop:s=1_unimodal}, the \emph{hook-multiplicity sequence} $m_{0,(n)},\,m_{1,(n)},\ldots, m_{n-1,(n)}$ is unimodal; 
We conjecture that it is unimodal for all partitions $\lambda\vdash n$.  

\begin{conjecture}\label{conj:unimodal}
    For every partition $\lambda \vdash n$, the sequence
    \[
        m_{0,\lambda},m_{1,\lambda},\ldots,m_{n-1,\lambda}
    \]
    is unimodal.
\end{conjecture}
    
The conjecture has been verified for all partitions of size $n \le 15$ and for all partitions of rectangular shape $(r^s)$ with $r \le 40$ and $s \le 5$.   
    
Note that, by Lemma~\ref{t:unimodal}, Conjecture~\ref{conj:unimodal} would provide an alternative proof of Theorem~\ref{t:main_hook-alternating2}. 
  
A sequence $a_0,\ldots,a_n$ of real numbers is 
\emph{log-concave} if  $a_{i-1}a_{i+1}\le a_i^2$ for all  $0<i<n$. 
It is not hard to show that a log-concave sequence with no internal zeros is unimodal, see e.g.~\cite{Brenti, Stanley}.
Since log-concavity implies unimodality, it is tempting to ask whether the hook-multiplicity sequence is log-concave.

\begin{conjecture}\label{conj:r1}
For every partition $\lambda=(r^s)$ with even $r\ne 6$,
the hook-multiplicity sequence $m_{0,\lambda},\,m_{1,\lambda},\ldots, m_{n-1,\lambda}$ is log-concave.
\end{conjecture}

Conjecture~\ref{conj:r1} 
was verified for all $r\leq 40$ and $s\leq 5$. 

Finally, recall that our proof of Theorem~\ref{thm:main} is not constructive. 
We conclude the paper with the following challenging 
problem.

\begin{problem} 
    Find 
    an explicit combinatorial description of a cyclic descent extension on the conjugacy class of each cycle type not equal to $(r^s)$ for a square-free $r$.
\end{problem}

A solution of this problem for the conjugacy classes of involutions 
is presented in~\cite{Adin-R23}. The analogous problem for standard Young tableaux of fixed non-ribbon shape was solved by Huang~\cite{Huang}. 

\longthanks{The second author is grateful for the hospitality of Bar-Ilan University during his Sabbatical leave.
Special thanks are due to Jan Saxl, whose personality influenced the development of this paper in many ways; in particular, the paper is inspired by the  
ingenious use of a higher Lie character in~\cite{Saxletal}.
Thanks also 
to Eli Bagno, Jonathan Bloom, Sergi Elizalde, Darij Grinberg, 
Vic Reiner, Richard Stanley, Sheila Sundaram and Josh Swanson for their comments.
}

